\chapterimage{figs/chapterhead.png}
\chapter{Application Overview}
\label{chap:application-overview}

This chapter maps the current GenESyS application family and explains which
executable a reader should open for each common workflow. It is intentionally
grounded in the current repository state, not in historical names or planned
future layouts.

The overview is split into four groups:
\begin{itemize}
\item the main interactive entry points;
\item the standalone Qt frontends;
\item the model-specific application route;
\item the planned-but-not-implemented DOE route.
\end{itemize}

Figure~\ref{fig:application-family-tree} summarizes that family tree. Chapter~\ref{chap:application-overview} then introduces the entry-point map in
Section~\ref{sec:application-overview-entry-points}, expands the core entry
points in Section~\ref{sec:application-overview-core-entry-points}, covers the
standalone Qt frontends in Section~\ref{sec:application-overview-standalone-qt-frontends},
describes the model-specific route in Section~\ref{sec:application-overview-model-specific},
and closes with the final choice flow in Section~\ref{sec:application-overview-choice}.
The following sections spell out the target, executable, purpose,
dependencies, public, status, evidence level, limitations, and the chapter
where each application is documented in more detail.

% FIGURE-SPEC-BEGIN
% ID: fig:application-family-tree
% Source of truth:
%   - source/applications/gui/CMakeLists.txt
%   - source/applications/gui/genesys/CMakeLists.txt
%   - source/applications/gui/httpworker/CMakeLists.txt
%   - source/applications/gui/dataanalyser/CMakeLists.txt
%   - source/applications/gui/optimizer/CMakeLists.txt
%   - source/applications/gui/ai_assistant/CMakeLists.txt
%   - source/applications/modelSpecific/CMakeLists.txt
%   - docs/ai_assistants/STATUS.md
% Purpose:
%   Show the current GenESyS application families and their relationships.
% Required visual content:
%   - main GUI as the primary entry point
%   - shell and worker as core support applications
%   - HTTP Worker GUI, Data Analyser, Optimizer, and AI Assistant as standalone Qt frontends
%   - model-specific applications as a selectable family
%   - Do Experiments as planned but not implemented
% Accessibility description:
%   A family-tree diagram that lets the reader choose the correct entry point quickly.
% Validation:
%   Each node must correspond to a real target, executable, or intentional absence in the current tree.
% FIGURE-SPEC-END
\begin{figure}[htbp]
\centering
\figureplaceholder{figs/generated/application-family-tree}
\caption{Current GenESyS application family and intended usage groups.}
\label{fig:application-family-tree}
\end{figure}

\section{Choosing the right entry point}
\label{sec:application-overview-entry-points}

If the goal is interactive model work, start with the main GUI. If the goal is
scripted or terminal-driven execution, start with the shell. If the goal is to
serve or control jobs through the worker runtime, use the worker. If the goal
is a specialized Qt helper, use one of the standalone frontends. If the goal
is to run a selected example from \texttt{source/applications/modelSpecific},
use the model-specific application route. If the goal is DOE functionality, no
application exists yet and the plan remains intentionally unimplemented.

Figure~\ref{fig:application-choice-flow} condenses that choice into a simple
decision path for readers who want the shortest route to the right executable.

\section{Core entry points}
\label{sec:application-overview-core-entry-points}

\subsection{Main GUI}
\label{subsec:application-overview-main-gui}

The main GUI is the canonical interactive entry point for model editing and
simulation work. It is the first application a new user should open when the
goal is to inspect, edit, or execute a model visually.

\noindent\textbf{Target.} \texttt{genesys\_gui\_application}, exposed through the aggregate \texttt{genesys\_gui} target.

\noindent\textbf{Executable.} \texttt{genesys-gui}.

\noindent\textbf{Purpose.} Primary Qt editor and simulation front end for day-to-day model work.

\noindent\textbf{Dependencies.} Qt6 Core, Gui, Widgets, PrintSupport, and Network, with Qt5 fallback in the current CMake code, plus the shared kernel, parser, plugin, simulator, worker-support, and tools libraries linked by the GUI target.

\noindent\textbf{Public.} General users, model authors, and reviewers who need the main interactive workspace.

\noindent\textbf{Status.} Partially validated.

\noindent\textbf{Evidence.} The installed launcher, build-tree startup, and About-dialog version check are documented in Chapter~\ref{chap:user-installation}; the current status ledger records focused GMDD diagnostics but not a full end-to-end interaction certification.

\noindent\textbf{Limitations.} Startup and focused GUI checks do not prove complete workflows or scientific correctness.

\noindent\textbf{Chapter.} Chapter~\ref{chap:main-graphical-application}.

\subsection{Shell}
\label{subsec:application-overview-shell}

The shell is the command-line entry point for scripted or terminal-based work.
It is the right choice when the user wants direct control over model loading,
commands, and repeatable text-driven execution.

\noindent\textbf{Target.} \texttt{genesys\_shell}.

\noindent\textbf{Executable.} \texttt{genesys\_shell}.

\noindent\textbf{Purpose.} Command-line front end for scripted model loading, command routing, and interactive terminal usage.

\noindent\textbf{Dependencies.} Kernel utility and statistics libraries, parser, plugins, and simulator support/runtime libraries.

\noindent\textbf{Public.} Advanced users, automation scripts, and developers who prefer a terminal workflow.

\noindent\textbf{Status.} Partially validated.

\noindent\textbf{Evidence.} The current status ledger records preset/build coverage, scripted command handling, plugin count reporting, and exit behavior.

\noindent\textbf{Limitations.} This chapter does not claim GUI workflows, file-based plugin autoload, or broader deployment behavior for the shell.

\noindent\textbf{Chapter.} Chapter~\ref{chap:shell-usage}.

\subsection{Worker}
\label{subsec:application-overview-worker}

The worker is the HTTP/background runtime that serves simulation-related
requests. It is not a GUI helper and it is not approved for public Internet
exposure by default.

\noindent\textbf{Target.} \texttt{genesys\_worker\_app}.

\noindent\textbf{Executable.} \texttt{genesys-worker}.

\noindent\textbf{Purpose.} Service runtime for controlled worker jobs, session handling, and backend request processing.

\noindent\textbf{Dependencies.} \texttt{genesys\_worker\_core} plus the shared kernel, parser, plugins, simulator support/runtime libraries, and the worker HTTP/session/auth stack.

\noindent\textbf{Public.} Operators and integration code in controlled academic or intranet-style deployments.

\noindent\textbf{Status.} Partially validated.

\noindent\textbf{Evidence.} The current status ledger records preset/build coverage, loopback health checks, exact JSON responses, and bounded exit behavior.

\noindent\textbf{Limitations.} Public bind, authentication, TLS, quotas, isolation, and audit controls are still separate security requirements, so startup evidence does not establish deployment readiness.

\noindent\textbf{Chapter.} Chapter~\ref{chap:worker-usage-and-security}.

\section{Standalone Qt frontends}
\label{sec:application-overview-standalone-qt-frontends}

Figure~\ref{fig:application-executable-matrix} groups the executable targets
with the shared backend libraries they rely on. This makes the current build
surface easier to scan without pretending that the helper windows are a single
monolithic executable.

% FIGURE-SPEC-BEGIN
% ID: fig:application-executable-matrix
% Source of truth:
%   - source/applications/gui/CMakeLists.txt
%   - source/applications/gui/genesys/CMakeLists.txt
%   - source/applications/gui/httpworker/CMakeLists.txt
%   - source/applications/gui/dataanalyser/CMakeLists.txt
%   - source/applications/gui/optimizer/CMakeLists.txt
%   - source/applications/gui/ai_assistant/CMakeLists.txt
% Purpose:
%   Show which GUI targets are paired with which shared backend libraries.
% Required visual content:
%   - one row per executable target
%   - the main GUI, HTTP Worker GUI, Data Analyser, Optimizer, and AI Assistant targets
%   - the shared runtime libraries used by each target
% Accessibility description:
%   A matrix that helps readers understand the separation between executables and reusable backends.
% Validation:
%   Each target and backend family must match the current CMake files.
% FIGURE-SPEC-END
\begin{figure}[htbp]
\centering
\figureplaceholder{figs/generated/application-executable-matrix}
\caption{Executable targets and shared backend library families for the current GUI applications.}
\label{fig:application-executable-matrix}
\end{figure}

\subsection{HTTP Worker GUI}
\label{subsec:application-overview-http-worker-gui}

The HTTP Worker GUI is the control dialog for the worker runtime. It exists to
support the worker-focused deployment path and to provide a graphical front end
around the worker service layer.

\noindent\textbf{Target.} \texttt{genesys\_httpworker\_gui\_application}.

\noindent\textbf{Executable.} \texttt{genesys-httpworker-gui}.

\noindent\textbf{Purpose.} Qt control window for worker-related configuration and interaction.

\noindent\textbf{Dependencies.} Qt6 or Qt5 Core, Gui, and Widgets, plus \texttt{genesys\_worker\_core}.

\noindent\textbf{Public.} Operators and developers who need a graphical worker control surface.

\noindent\textbf{Status.} Defined in CMake, but not independently validated beyond the current target and preset evidence.

\noindent\textbf{Evidence.} The current build files and presets define the application route and executable name.

\noindent\textbf{Limitations.} There is no separate startup or interaction certification yet, so this application should be treated as present but lightly evidenced.

\noindent\textbf{Chapter.} Chapter~\ref{chap:standalone-tools}.

\subsection{Data Analyser}
\label{subsec:application-overview-data-analyser}

The Data Analyser is a specialized Qt window for data-analysis workflows. It
is useful when the user wants a dedicated analysis surface instead of the main
model editor.

\noindent\textbf{Target.} \texttt{genesys\_dataanalyser\_gui\_application}.

\noindent\textbf{Executable.} \texttt{genesys-dataanalyser-gui}.

\noindent\textbf{Purpose.} Dedicated analysis frontend for imported or prepared data.

\noindent\textbf{Dependencies.} Qt6 or Qt5 Core, Gui, and Widgets, plus the shared kernel, parser, plugins, simulator support/runtime, and tools libraries.

\noindent\textbf{Public.} Users who need a dedicated analysis workspace.

\noindent\textbf{Status.} Startup validated.

\noindent\textbf{Evidence.} The current status ledger records Xvfb window creation, liveness, and controlled teardown.

\noindent\textbf{Limitations.} Startup does not establish analysis correctness, workflow completeness, or Level 3 scientific maturity.

\noindent\textbf{Chapter.} Chapter~\ref{chap:standalone-tools}.

\subsection{Optimizer}
\label{subsec:application-overview-optimizer}

The Optimizer is the optimization-oriented Qt frontend. It is the GUI entry
point for optimization work when the user wants a standalone workspace instead
of the main editor.

\noindent\textbf{Target.} \texttt{genesys\_optimizer\_gui\_application}.

\noindent\textbf{Executable.} \texttt{genesys-optimizer-gui}.

\noindent\textbf{Purpose.} Dedicated optimization frontend for the current optimizer scaffold.

\noindent\textbf{Dependencies.} Qt6 or Qt5 Core, Gui, and Widgets, plus the shared kernel, parser, plugins, simulator support/runtime, and tools libraries.

\noindent\textbf{Public.} Advanced users and developers experimenting with optimization workflows.

\noindent\textbf{Status.} Startup validated.

\noindent\textbf{Evidence.} The current status ledger records Xvfb window creation, liveness, and controlled teardown.

\noindent\textbf{Limitations.} Startup does not prove a selected optimization algorithm, full workflow maturity, or validated scientific output.

\noindent\textbf{Chapter.} Chapter~\ref{chap:standalone-tools}.

\subsection{AI Assistant}
\label{subsec:application-overview-ai-assistant}

The AI Assistant is the dedicated Qt frontend for assistant-oriented
interaction. It is separated from the main GUI so the assistant workflow can be
handled independently.

\noindent\textbf{Target.} \texttt{genesys\_ai\_assistant\_gui\_application}.

\noindent\textbf{Executable.} \texttt{genesys-ai-assistant-gui}.

\noindent\textbf{Purpose.} Dedicated assistant UI for prompt, tool, and response workflows.

\noindent\textbf{Dependencies.} Qt6 or Qt5 Core, Gui, and Widgets, plus the shared kernel, parser, plugins, simulator support/runtime, and tools libraries.

\noindent\textbf{Public.} Users and developers who need a separate assistant-oriented application surface.

\noindent\textbf{Status.} Startup validated.

\noindent\textbf{Evidence.} The current status ledger records no-credential Xvfb startup and teardown.

\noindent\textbf{Limitations.} Startup does not establish provider integration, credential handling, redaction, or failure workflow maturity.

\noindent\textbf{Chapter.} Chapter~\ref{chap:standalone-tools}.

\section{Model-specific applications}
\label{sec:application-overview-model-specific}

Figure~\ref{fig:application-maturity-map} places the current families on a
simple maturity ladder. The intent is to prevent overclaiming: different
applications are at different evidence levels, and startup alone is not the
same as full workflow validation.

% FIGURE-SPEC-BEGIN
% ID: fig:application-maturity-map
% Source of truth:
%   - docs/ai_assistants/STATUS.md
%   - docs/ai_assistants/ARCHITECTURE.md
%   - docs/ai_assistants/reference/APPLICATIONS_TOOLS_MODELS.md
% Purpose:
%   Show the current maturity and evidence level of each application family.
% Required visual content:
%   - partially validated core applications
%   - startup-validated standalone Qt frontends
%   - snapshot-only model-specific applications
%   - intentionally absent Do Experiments route
% Accessibility description:
%   A maturity chart that distinguishes availability from validation depth.
% Validation:
%   The labels must match the current status ledger and architecture notes.
% FIGURE-SPEC-END
\begin{figure}[htbp]
\centering
\figureplaceholder{figs/generated/application-maturity-map}
\caption{Current application maturity and evidence levels.}
\label{fig:application-maturity-map}
\end{figure}

The model-specific application route builds one selected example from
\path{source/applications/modelSpecific}. The current CMake logic treats the
selection as a build-time input, so the exact executable contents depend on the
selected source file.

\noindent\textbf{Target.} \path{genesys_modelspecific_app}. The selected source file is provided by \path{GENESYS_MODELSPECIFIC_APP}.

\noindent\textbf{Executable.} \path{genesys_modelspecific_app}. A legacy compatibility copy named \path{genesys_terminal_application} is created after the build.

\noindent\textbf{Purpose.} Build and run one selected model-specific example or teaching application from the model-specific tree.

\noindent\textbf{Dependencies.} The shared kernel, parser, plugins, simulator support/runtime libraries, and the selected model-specific source and header files.

\noindent\textbf{Public.} Advanced users, instructors, and maintainers who need concrete example flows.

\noindent\textbf{Status.} Snapshot only.

\noindent\textbf{Evidence.} The current status ledger treats the family as a historical bounded sweep that still needs current revalidation.

\noindent\textbf{Limitations.} There is no single fixed executable payload; the selected example determines the runtime behavior, and each selection must be rechecked on the current branch.

\noindent\textbf{Chapter.} Chapter~\ref{chap:model-specific-and-advanced-examples}.

\subsection{Do Experiments}
\label{subsec:application-overview-do-experiments}

The DOE route is intentionally not implemented yet. The CMake umbrella knows
that the feature is planned, but the repository does not provide a real empty
application for it.

\noindent\textbf{Target.} None yet.

\noindent\textbf{Executable.} None yet.

\noindent\textbf{Purpose.} Future design-of-experiments functionality, once a bounded backend and product contract exist.

\noindent\textbf{Dependencies.} No implemented application route yet.

\noindent\textbf{Public.} None yet.

\noindent\textbf{Status.} Not started; intentionally unavailable.

\noindent\textbf{Evidence.} The GUI umbrella emits a fatal configuration error when the feature flag is enabled, which confirms that the repository does not ship a placeholder executable.

\noindent\textbf{Limitations.} Readers should not expect a launcher, preset, or runnable window for this route until the backend/product definition is completed.

\noindent\textbf{Chapter.} No dedicated user chapter yet.

\section{How to choose}
\label{sec:application-overview-choice}

Figure~\ref{fig:application-choice-flow} turns the application map into a
decision tree. The reader should use it as a quick answer to the question
``which application should I open first?'' without needing to inspect the full
source tree. For direct lookup, Section~\ref{subsec:application-overview-main-gui}
covers the main GUI, Section~\ref{subsec:application-overview-shell} covers the
shell, Section~\ref{subsec:application-overview-worker} covers the worker,
Section~\ref{subsec:application-overview-http-worker-gui} covers the HTTP
Worker GUI, Section~\ref{subsec:application-overview-data-analyser} covers the
Data Analyser, Section~\ref{subsec:application-overview-optimizer} covers the
Optimizer, Section~\ref{subsec:application-overview-ai-assistant} covers the
AI Assistant, Section~\ref{sec:application-overview-model-specific} covers
the model-specific route, and Section~\ref{subsec:application-overview-do-experiments}
covers the intentionally absent DOE path.

% FIGURE-SPEC-BEGIN
% ID: fig:application-choice-flow
% Source of truth:
%   - docs/ai_assistants/STATUS.md
%   - docs/ai_assistants/ARCHITECTURE.md
%   - source/applications/gui/CMakeLists.txt
%   - source/applications/modelSpecific/CMakeLists.txt
% Purpose:
%   Show how a reader chooses the correct GenESyS application for a task.
% Required visual content:
%   - decision path for main GUI, shell, worker, standalone tools, model-specific apps, and no-DOE path
%   - a clear branch for intentionally absent Do Experiments support
% Accessibility description:
%   A simple flowchart that avoids overloading the reader with build details.
% Validation:
%   The flow must match the current application inventory and status notes.
% FIGURE-SPEC-END
\begin{figure}[htbp]
\centering
\figureplaceholder{figs/generated/application-choice-flow}
\caption{Decision flow for selecting the right GenESyS application.}
\label{fig:application-choice-flow}
\end{figure}

In practical terms, the choice is simple:
\begin{itemize}
\item edit or inspect a model visually: open the main GUI;
\item run scripts or terminal commands: open the shell;
\item serve controlled jobs: open the worker;
\item use a helper window for analysis, optimization, or assistant tasks: open the matching standalone frontend;
\item run a selected teaching or example model: use the model-specific application route;
\item look for DOE: do not expect an application yet.
\end{itemize}

The rest of the user manual expands each of those choices in the dedicated
chapters listed above.
